%%%%%%%%%%%%%%%%%%%%%%%%%%%%%%%%%%%%%%%%%%%%%%%%%
% Packages %
%%%%%%%%%%%%%%%%%%%%%%%%%%%%%%%%%%%%%%%%%%%%%%%%%

%Usual ones
\usepackage[french]{babel}
\usepackage[T1]{fontenc}
\usepackage{lmodern}
\usepackage[utf8]{inputenc}
\usepackage{csquotes} %Guillemets français
\usepackage{amsmath}
\usepackage{amssymb}
\usepackage{amsfonts}
\usepackage{braket}
\usepackage{multirow}
\usepackage{wasysym} %Pour des intégrales doubles fermées
\usepackage{color}
\usepackage{slashed}
\usepackage{enumerate}
\usepackage{enumitem}
\usepackage[output-open-uncertainty=\pm,output-close-uncertainty=,
output-decimal-marker={,}, inter-unit-product={}\cdot{}, range-units = single, range-phrase=\text{ à },separate-uncertainty=true]
{siunitx}
\sisetup{detect-all}
\usepackage{todonotes}
\usepackage{epigraph}
\usepackage{tabularx}
\usepackage{longtable}
\usepackage{fontawesome}
\usepackage[tikz]{bclogo} %Helps for creation of colored boxes
\usepackage{caption}
%\usepackage{graphbox}
\usepackage{listingsutf8}
\usepackage{pythonhighlight} %Pour insérer du python
%\usepackage{mhchem}
%%%%%%%%%%%%%%%%%%%%%%

%Packages used to provide commands depending on an argument
\usepackage{ifthen}
\usepackage{ifmtarg}
\usepackage{etoolbox}

%%%%%%%%%%%%%%%%%%%%%%


%tikz
\usepackage{tikz}
\usepackage{tikz-3dplot}
\usepackage[europeanresistor]{circuitikz}
\usetikzlibrary{%
	shapes.geometric,calc,intersections,%
	decorations.markings,decorations.text,decorations,%
	patterns,decorations.pathmorphing,fpu,%
	positioning,automata,shapes.multipart,%
	arrows.meta, backgrounds,%
	tikzmark, babel, perspective
}


%%%%%%%%%%%%%%%%%%%%%%%
% Page geometry
\usepackage[top=2cm, bottom=2cm, left=1.5cm, right=1.5cm]{geometry}
	
%%%%%%%%%%%%%%%%%%%%%%

%Appendices
\usepackage[toc,page]{appendix}

%%%%%%%%%%%%%%%%%%%%%%

%Bibliography
% \usepackage[square]{natbib}
\usepackage[nottoc]{tocbibind} 
\usepackage[
	backend=biber,
	defernumbers=true,
	natbib=true,
	maxbibnames=99,
	giveninits=true,
	style=numeric,
	url=false,
	isbn=false]{biblatex}
\addbibresource{bib/biblio.bib}


%%%%%%%%%%%%%%%%%%%%%%%%%

%hyperref
\usepackage[%
  pdftex,colorlinks=true,
	pdfstartview=FitV,
	linkcolor= linkcolor,
	citecolor= linkcolor,
	urlcolor= linkcolor,
	hyperindex=true,
	hyperfigures=true]{hyperref}
	
%%%%%%%%%%%%%%%%%%%%%%%%

%Clever Referencing
\usepackage[noabbrev]{cleveref}
\crefname{appendixchapter}{annexe}{annexes}
\Crefname{appendixchapter}{Annexe}{Annexes}
\crefname{appendixsection}{section annexe}{Sections annexes}
\Crefname{appendixsection}{Section annexe}{Sections annexes}
	
%%%%%%%%%%%%%%%%%%%%%%%%%%%%%%%%%%%%%%%%%%%%%%%%%%%%%%%%%
% Defined commands and functions
%%%%%%%%%%%%%%%%%%%%%%%%%%%%%%%%%%%%%%%%%%%%%%%%%%%%%%%%%

%Linkcolor
\definecolor{linkcolor}{rgb}{0.7,0.2,0.2}

%%%%%%%%%%%%%%%%%%%%%%%%%%

%Short usual commands

\newcommand{\me}{\mathrm{e}}
\newcommand{\mi}{\mathrm{i}}
\newcommand{\mj}{\mathrm{j}}
\newcommand{\md}{\mathrm{d}}
\newcommand{\eo}{\epsilon_0}
\newcommand{\heaviside}{\mathrm{H}}
\newcommand{\ve}[1]{\vbox{\halign{##\cr 
  \tiny\rightarrowfill\cr\noalign{\nointerlineskip\vskip1pt} 
  $#1\mskip2mu$\cr}}}
\newcommand{\co}[1]{\underline{#1}}
\newcommand{\tf}[1]{\hat{#1}}

%%%%%%%%%%%%%%%%%%%%%%%%%%

%Mathematical operators

\DeclareMathOperator{\mRe}{Re}
\DeclareMathOperator{\mIm}{Im}
\DeclareMathOperator{\sinc}{sinc}
\DeclareMathOperator{\trace}{Tr}
\DeclareMathOperator{\diverg}{div}
\DeclareMathOperator{\grad}{\ve{\mathrm{grad}}}
\DeclareMathOperator{\rot}{\ve{\mathrm{rot}}}
\DeclareMathOperator{\laplac}{\Delta}

%%%%%%%%%%%%%%%%%%%%%%%%%%%

%Special trick to allow moving of the \in character up or down, by 
%putting a length as first argument
%Usage for sum indices : n \in[-0.6pt] \mathbb{R} for example
\newcommand{\myin}[1][0pt]{%
  \mathchoice{\raisebox{#1}{$\displaystyle\in$}}
             {\raisebox{#1}{$\in$}}
             {\raisebox{#1}{$\scriptstyle\in$}}
             {\raisebox{#1}{$\scriptscriptstyle\in$}}
             }
\makeatletter

%%%%%%%%%%%%%%%%%%%%%%%%%%

%Personal todo command
\newcommand{\mytodo}[1]{{\todo[inline,backgroundcolor=blue!10!white]
{#1}}}

%%%%%%%%%%%%%%%%%%%%%%%%%%

%Defining commands allowing to create parts that are printed
%only for the teacher and/or only for students.
%IMPORTANT : the forstudent toggle MUST be set to True or False
%when compiling the document

\newtoggle{forstudent}

%This command allows to create teacher notes, that should not be
%visible to students. 
\newcommand{\teacher}[1]{%
	\iftoggle{forstudent}%
		{}	%nothing printed if forstudent=true
		{#1}%print the text if forstudent=false
	}
	
%This command allows to create student notes, that should be printed
%only in the student version
\newcommand{\student}[1]{%
	\iftoggle{forstudent}%
		{#1}%print the text if forstudent=true
		{}	%nothing printed if forstudent=false
	}

%%%%%%%%%%%%%%%%%%%%%%%%%%

%Defining some specific useful environments

%This command creates a small box in which the text will be displayed. 
\newcommand{\important}[1]{\begin{center}\fbox{
\vspace{0.1cm}
\begin{minipage}{0.9\textwidth}
#1
\end{minipage}
\vspace{0.1cm}
}\end{center}
}

%This command can be used for exemples
\newcommand{\exemple}[1]{\begin{center}
\vspace{0.1cm}
\begin{minipage}{0.9\textwidth}
\textbf{Exemple} :
#1
\end{minipage}
\vspace{0.1cm}
\end{center}
}

%This command creates a specific environment for oral information
%to be given. Can only be displayed in the teacher version.
\newcommand{\oral}[1]{%
	\teacher{	
		\begin{bclogo}[avecBarre = false,arrondi = 0.1,
		couleur = green!20!white,logo=\ ]{}
		\vspace{-0.3cm} \faComment~#1  
		\end{bclogo}
		}
	}
	
%This command creates a box for remarks.
\newcommand{\remark}[1]{%
	\begin{bclogo}[avecBarre = true,noborder=true,logo=\ ]{}
		\textcolor{gray}{
	 #1  }
	\end{bclogo}
	}
	